\chapter{Expectations and the Lucas Critique}\label{ch:b2:expectations-and-lucas-critique}

How people form expectations decides what policy can do. This chapter takes expectations seriously: it asks what
happens to the relations economists estimate when people understand the policy that is being run, and why a central bank
that can always choose its policy afresh may end up with more inflation than it wants. The chart sets average inflation in
fifteen rich countries beside a measure of how independent their central banks were by law.

\begin{openingchart}{Average annual inflation 1955--1988 against an index of the legal independence of the central bank
  (1 least, 4 most independent), fifteen countries; the line is the least-squares fit.}{Jord\`a--Schularick--Taylor
  Macrohistory Database (cpi); United States: Bureau of Labor Statistics, CPI-U; Alesina and Summers
  (1993), Table A1; our computation}
\begin{tikzpicture}
\begin{axis}[omrecord, width=0.62\linewidth, height=0.44\linewidth, xmin=1, xmax=4.25, ymin=2, ymax=11,
    xlabel={Index of central-bank independence}, ylabel={Average inflation, per cent},
    nodes near coords, point meta=explicit symbolic,
    nodes near coords style={font=\scriptsize, anchor=west, xshift=1pt}]
  \addplot[only marks, mark=*, omDef, mark size=1.6pt, nodes near coords={}] table[x=index, y=infl]
    {figdata/bachelor-2/25-expectations-and-lucas-critique/out/cbi.csv};
  \addplot[only marks, mark=none] table[x=index, y=infl, meta=iso3]
    {figdata/bachelor-2/25-expectations-and-lucas-critique/out/cbi_lab.csv};
  \node[font=\scriptsize, anchor=east, omDef, align=right] at (axis cs:1.96,5.8) {AUS, FRA, GBR,\\NOR, SWE};
  \node[font=\scriptsize, anchor=east, omDef] at (axis cs:2.46,4.98) {CAN, JPN};
  \node[font=\scriptsize, anchor=north east, omDef] at (axis cs:4.1,3.05) {CHE, DEU};
  \addplot[omThm, thick, dashed, domain=1:4.25, samples=2, nodes near coords={}] {10.29 - 1.83*x};
\end{axis}
\end{tikzpicture}
\end{openingchart}

% figdata: cbi (ESP 9.95, DEU 3.19, CHE 3.42), cbi_fit (slope -1.83, intercept 10.29, corr -0.80, n 15)
Spain, with the least independent central bank of the group, averaged 9.95 per cent inflation over the period; Germany and
Switzerland, with the most independent, 3.19 and 3.42 per cent. Each point of the index goes with 1.83 points less inflation,
and the correlation is $-0.80$. Several countries share an index value, six of them at 2, and their labels are grouped. A model of
why independence might lower inflation is in \cref{sec:b2:expectations-and-lucas-critique:time}; whether it does is the
debate of \cref{sec:b2:expectations-and-lucas-critique:credibility}.

\begin{recall}
From the Year 1 volume: adaptive expectations, the expectations-augmented Phillips curve, policy rules and discretion.
From the school volumes: the Phillips curve and stagflation. From \cref{ch:b2:ramsey-model}: households that look ahead.
\end{recall}

\begin{notation}[Symbols of this chapter]\label{not:b2:expectations-and-lucas-critique:symbols}
$\pi$ is inflation and $\pi^e$ expected inflation, set before $\pi$ is known; $u$ is unemployment and $u^*$ its natural rate; $a > 0$ is the effect of an inflation
surprise on unemployment; $e$ is a supply shock with mean zero and variance $\sigma^2$; $\pi^{\mathrm T}$ is the inflation target; $k > 0$ is how far below $u^*$ policy would like
unemployment to be; $\lambda$ is society's weight on unemployment, $w$ the central bank's. $E_{t-1}x$ is the expectation of $x$ formed with the
information available at $t - 1$.
\end{notation}

\section{Adaptive and rational expectations}\label{sec:b2:expectations-and-lucas-critique:rational}

Under adaptive expectations, $\pi^e_t = \pi_{t-1}$ or a weighted average of past inflation. People who form expectations this way make
mistakes that are predictable: when inflation rises steadily, they underpredict it year after year, and a policy that
raises inflation steadily keeps unemployment below its natural rate for as long as it lasts.

\begin{definition}[Rational expectations]\label{def:b2:expectations-and-lucas-critique:rational-expectations}
Expectations are \emph{rational}\index{rational expectations} when they equal the mathematical expectation of the
variable implied by the model itself, given the information available: $\pi^e_t = E_{t-1}\pi_t$.
\end{definition}

The forecast error $\pi_t - E_{t-1}\pi_t$ then has mean zero and is uncorrelated with anything known at $t - 1$: if it were correlated with something
known, the forecast could be improved by using it. \omterm{def:b2:expectations-and-lucas-critique:rational-expectations}{Rational expectations} do not mean that people know the future; they mean
that people do not keep making the same mistake. In particular, people understand the policy rule being followed, and its
systematic part is built into their expectations.

\section{The Lucas critique}\label{sec:b2:expectations-and-lucas-critique:critique}

\begin{definition}[Lucas critique]\label{def:b2:expectations-and-lucas-critique:lucas-critique}
The \emph{Lucas critique}\index{Lucas critique} is the argument that the parameters of an estimated relation between
economic variables can change when policy changes, because they depend on the expectations people form about that policy,
so the relation cannot be used to predict the effects of a new policy.
\end{definition}

Take the Phillips curve $u_t = u^* - a(\pi_t - \pi^e_t)$, and suppose the central bank keeps inflation at a target $\pi^{\mathrm T}$ plus a random policy surprise $\varepsilon_t$ that no one can forecast.
With \omterm{def:b2:expectations-and-lucas-critique:rational-expectations}{rational expectations} $\pi^e_t = \pi^{\mathrm T}$, so $u_t = u^* - a(\pi_t - \pi^{\mathrm T}) = (u^* + a\pi^{\mathrm T}) - a\pi_t$. An economist who regresses unemployment on inflation in this regime finds a stable
negative relation with intercept $u^* + a\pi^{\mathrm T}$ and slope $-a$. But the intercept contains the policy: if the bank raises the target, the curve shifts
up by $a$ times the rise and unemployment, on average, stays at $u^*$. The listing estimates the curve under a 2 per cent target and
then raises the target to 4.

\omcode{code/bachelor-2/25-expectations-and-lucas-critique/listings/lucas.py}{1}{29}{An estimated Phillips curve used to
predict a change of policy.}\label{lst:b2:expectations-and-lucas-critique:lucas}

\omoutput{code/bachelor-2/25-expectations-and-lucas-critique/outputs/lucas.txt}{The estimated curve promises
unemployment of about 1 per cent at 4 per cent inflation; the economy delivers its natural rate, and the curve moves.}\label{out:b2:expectations-and-lucas-critique:lucas}

% code: listings/lucas.py
The regression on the first forty years is excellent: slope $-1.91$, close to the true $-2$. Read as a menu, it says that 4 per cent
inflation would bring unemployment to 1.27 per cent. Under the new target inflation averages 3.81 per cent and unemployment 5.38 per cent,
the natural rate of 5 plus noise; the re-estimated curve has much the same slope and an intercept higher by about 4.3, close to
$a \times 2 = 4$. Nothing in the first regime's data could have revealed this, because the target never changed in it.

\begin{omchart}{United States, inflation against unemployment, annual averages of quarterly data, 1960--2008, by decade.}{statsmodels
  macrodata (US national accounts and labour data, quarterly 1959--2009); our computation}
\begin{tikzpicture}
\begin{axis}[omrecord, width=0.62\linewidth, height=0.44\linewidth, xmin=3, xmax=10, ymin=-1, ymax=12,
    xlabel={Unemployment rate, per cent}, ylabel={Inflation, per cent},
    legend style={at={(1.03,0.5)}, anchor=west, font=\scriptsize}]
  \addplot[only marks, mark=*, omDef, mark size=1.6pt] table[x=unemp, y=infl, restrict expr to domain={\thisrow{decade}}{1960:1960}]
    {figdata/bachelor-2/25-expectations-and-lucas-critique/out/phillips.csv};
  \addlegendentry{1960s}
  \addplot[only marks, mark=square*, omThm, mark size=1.6pt] table[x=unemp, y=infl, restrict expr to domain={\thisrow{decade}}{1970:1970}]
    {figdata/bachelor-2/25-expectations-and-lucas-critique/out/phillips.csv};
  \addlegendentry{1970s}
  \addplot[only marks, mark=triangle*, omMeth, mark size=2pt] table[x=unemp, y=infl, restrict expr to domain={\thisrow{decade}}{1980:1980}]
    {figdata/bachelor-2/25-expectations-and-lucas-critique/out/phillips.csv};
  \addlegendentry{1980s}
  \addplot[only marks, mark=o, black!70, mark size=1.6pt] table[x=unemp, y=infl, restrict expr to domain={\thisrow{decade}}{1990:2000}]
    {figdata/bachelor-2/25-expectations-and-lucas-critique/out/phillips.csv};
  \addlegendentry{1990--2008}
\end{axis}
\end{tikzpicture}
\end{omchart}

\begin{record}[The shifting Phillips curve, 1960--1989]
% figdata: phillips_dec (1960s 2.55 / 4.78; 1970s 7.22 / 6.22; 1980s 4.91 / 7.28)
In the 1960s American inflation and unemployment traced a downward-sloping curve: inflation averaged 2.55 per cent and
unemployment 4.78 per cent. In the 1970s both were higher on average, 7.22 and 6.22 per cent; in the 1980s inflation fell back
to 4.91 per cent while unemployment averaged 7.28. In the 1960s and the 1970s the points within the decade slope downward; in the 1980s they barely slope. Across decades the curve
moved up and then back, as expected inflation rose and then fell. A curve estimated on the 1960s promised low unemployment
in exchange for inflation; the 1970s delivered the inflation without the low unemployment.
\end{record}

\section{Policy ineffectiveness and its limits}\label{sec:b2:expectations-and-lucas-critique:pip}

In the Lucas supply model, output and unemployment depart from their natural levels only through surprises:
$u_t = u^* - a(\pi_t - E_{t-1}\pi_t)$. Suppose the bank sets inflation by a rule that responds to last year's unemployment, $\pi_t = \pi^{\mathrm T} - b(u_{t-1} - u^*) + \varepsilon_t$. Then $E_{t-1}\pi_t = \pi^{\mathrm T} - b(u_{t-1} - u^*)$,
the surprise is $\varepsilon_t$, and $u_t = u^* - a\varepsilon_t$ whatever the value of $b$.

\begin{definition}[Policy-ineffectiveness proposition]\label{def:b2:expectations-and-lucas-critique:pip}
The \emph{policy-ineffectiveness proposition}\index{policy-ineffectiveness proposition} states that, with \omterm{def:b2:expectations-and-lucas-critique:rational-expectations}{rational
expectations} and fully flexible prices, a systematic monetary policy rule cannot affect the behaviour of output or
unemployment; only unanticipated policy can.
\end{definition}

\begin{model}{Lucas supply model}
\assumes \omterm{def:b2:expectations-and-lucas-critique:rational-expectations}{Rational expectations}; prices and wages that are set each period with all the information available; output that
departs from its natural level only through price surprises.
\predicts Anticipated changes in money or inflation leave output unchanged; only surprises move it; the Phillips curve
shifts when the policy regime changes.
\tested Against the 1970s, when inflation rose without a lasting fall in unemployment: the shifting curve of the record.
\breaks Wages and prices are often set for a year or more in advance. A contract signed before a shock cannot respond to it,
so even a fully anticipated policy response, made after the contract, changes real outcomes: with sticky prices, systematic
policy matters.
\end{model}

The limit is the timing. If wages for year $t$ are fixed at $t - 2$, a rule that responds to the shock at $t - 1$ moves inflation after the
wage is fixed and before it can change; the response is anticipated by everyone, but too late for the contract. Models with
staggered contracts and sticky prices keep the \omterm{def:b2:expectations-and-lucas-critique:rational-expectations}{rational expectations} and drop the flexible prices.

\begin{omfigure}
\begin{tikzpicture}[econ]
  \draw[->, thick] (0,0) -- (12.6,0) node[right] {time};
  \foreach \x/\a/\b in {1/{Expectations formed,}/{contracts signed}, 4.4/{Shock}/{occurs}, 7.8/{Policy}/{responds}, 11.2/{Inflation and}/{unemployment}} {
    \draw[thick] (\x,-0.15) -- (\x,0.15);
    \node[align=center, font=\small, above] at (\x,0.25) {\a\\\b};
  }
  \node[font=\small, below] at (1,-0.2) {$\pi^e_t = E_{t-1}\pi_t$};
  \node[font=\small, below] at (4.4,-0.2) {$e_t$};
  \node[font=\small, below] at (7.8,-0.2) {$\pi_t$};
  \node[font=\small, below] at (11.2,-0.2) {$u_t$};
\end{tikzpicture}

{\small The order of events in every model of this chapter. Because expectations are formed before the shock, policy can
respond to the shock even though everyone knows the rule; it cannot surprise anyone systematically.}
\end{omfigure}

\section{Time inconsistency}\label{sec:b2:expectations-and-lucas-critique:time}

\begin{definition}[Central-bank loss function]\label{def:b2:expectations-and-lucas-critique:loss}
A \emph{central-bank loss function}\index{central-bank loss function} is a function that measures how bad an outcome is
for the central bank, here $L = \tfrac12(\pi - \pi^{\mathrm T})^2 + \tfrac12 w(u - u^* + k)^2$; the bank chooses policy to make the expected loss as small as
possible.
\end{definition}

The term $k > 0$ says that the bank would like unemployment below its natural rate, for instance because taxes or market
power make the natural rate too high. Society's own loss has the same form, with weight $\lambda$; for now $w = \lambda$.

\begin{definition}[Time inconsistency]\label{def:b2:expectations-and-lucas-critique:time-inconsistency}
A policy plan is \emph{time inconsistent}\index{time inconsistency} when the policy maker, free to choose again at a later
date, would choose differently from the plan even though nothing unforeseen has happened.
\end{definition}

Without shocks, $u = u^* - a(\pi - \pi^e)$. Under discretion the bank takes $\pi^e$ as given, since expectations are already set, and chooses $\pi$. The first-order
condition is $(\pi - \pi^{\mathrm T}) - wa(u - u^* + k) = 0$, that is
\begin{equation}\label{eq:b2:expectations-and-lucas-critique:reaction}
\pi = \frac{\pi^{\mathrm T} + wa^2\pi^e + wak}{1 + wa^2}.
\end{equation}
People know this. With \omterm{def:b2:expectations-and-lucas-critique:rational-expectations}{rational expectations} $\pi^e = \pi$, which gives $\pi - \pi^{\mathrm T} = wak$, and then $u = u^*$.

\begin{definition}[Inflation bias]\label{def:b2:expectations-and-lucas-critique:inflation-bias}
The \emph{inflation bias}\index{inflation bias} of discretionary policy is the excess of average inflation over the
target that results when the policy maker cannot commit, here $wak$.
\end{definition}

If the bank could commit to $\pi = \pi^{\mathrm T}$, expectations would be $\pi^{\mathrm T}$ and unemployment $u^*$, the same as under discretion, with less inflation. The commitment is \omterm{def:b2:expectations-and-lucas-critique:time-inconsistency}{time inconsistent}:
once expectations equal $\pi^{\mathrm T}$, the bank gains by a surprise, so an announcement of $\pi^{\mathrm T}$ is not believed. With $\pi^{\mathrm T} = 2$, $a = 0.5$, $k = 2$ and $w = 1$, discretion gives 3 per
cent inflation and a loss of $\tfrac12 \times 1 + \tfrac12 \times 4 = 2.5$; commitment gives 2 per cent and a loss of 2.

\begin{omfigure}
\begin{tikzpicture}[econ]
\begin{axis}[omrecord, grid=none, xmin=0.5, xmax=8.6, ymin=-0.5, ymax=5.5, axis equal image, width=0.7\linewidth,
    xlabel={Unemployment, $u$}, ylabel={Inflation, $\pi$}, xtick={3,5}, xticklabels={$u^* - k$, $u^*$},
    ytick={2,3}, yticklabels={$\pi^{\mathrm T}$, $\pi^{\mathrm T} + wak$}]
  \draw[black!40] (axis cs:5,-0.5) -- (axis cs:5,5.5);
  \addplot[omDef, thick, domain=3.8:6.3, samples=2] {3 - 2*(x - 5)};
  \addplot[omDef, thick, dashed, domain=3.3:5.8, samples=2] {2 - 2*(x - 5)};
  \addplot[omThm, thick, domain=0:360, samples=90, variable=\t] ({3 + sqrt(5)*cos(t)}, {2 + sqrt(5)*sin(t)});
  \addplot[omThm, thick, dashed, domain=0:360, samples=90, variable=\t] ({3 + 2*cos(t)}, {2 + 2*sin(t)});
  \node[eqpoint, label={[font=\small]right:D}] at (axis cs:5,3) {};
  \node[eqpoint, label={[font=\small]below left:C}] at (axis cs:5,2) {};
  \node[font=\scriptsize] at (axis cs:3,2) {$+$};
  \node[font=\small, anchor=west] at (axis cs:6.3,0.4) {$\pi^e = \pi^{\mathrm T} + wak$};
  \node[font=\small, anchor=east] at (axis cs:3.3,5.4) {$\pi^e = \pi^{\mathrm T}$};
\end{axis}
\end{tikzpicture}

{\small The \omterm{def:b2:expectations-and-lucas-critique:inflation-bias}{inflation bias} with $w = 1$, $a = 0.5$, $k = 2$. Each Phillips line holds for one level of expected inflation; the circles are loss contours around the bank's ideal point (cross).
Under discretion the bank chooses the point of tangency on the line it faces; only at D does that point also confirm the
expectations. Commitment C has the same unemployment, less inflation and a smaller loss contour, but on the dashed line the
bank would rather move along it to the north-west.}
\end{omfigure}

\begin{model}{The inflation-bias model}
\assumes \omterm{def:b2:expectations-and-lucas-critique:rational-expectations}{Rational expectations}; a Phillips curve with surprises; a policy maker who wants unemployment below its natural
rate and chooses inflation after expectations are set.
\predicts Average inflation above target by $wak$, with no gain in unemployment; lower inflation where the policy maker can commit or
cares less about unemployment.
\tested Across rich countries, more legally independent central banks went with lower average inflation, 1955--1988 (the
opening chart; Alesina and Summers 1993).
\breaks Some central banks with little legal independence kept inflation low, and fundamentals such as openness explain
much of the cross-country variation in inflation (Campillo and Miron 1997); a correlation leaves open whether
independence lowers inflation or whether countries averse to inflation choose independent banks.
\end{model}

\section{Credibility, independence and inflation targets}\label{sec:b2:expectations-and-lucas-critique:credibility}

\begin{definition}[Credibility]\label{def:b2:expectations-and-lucas-critique:credibility}
A policy is \emph{credible}\index{credibility (of policy)} when the public expects it to be carried out, because the
policy maker will want to carry it out when the time comes or cannot do otherwise.
\end{definition}

Three ways out of the \omterm{def:b2:expectations-and-lucas-critique:inflation-bias}{inflation bias} have been tried. The first is a \omterm{def:b2:repeated-games:reputation}{reputation}: a bank that faces the same public every
year may find that the gain from a surprise today is less than the cost of higher expected inflation for years after. The
second is delegation.

\begin{definition}[Central-bank independence]\label{def:b2:expectations-and-lucas-critique:cbi}
\emph{Central-bank independence}\index{central-bank independence} is the extent to which a central bank sets monetary
policy free of instructions from the government, by law and in practice.
\end{definition}

Delegating policy to a bank that weighs unemployment less than society does, $w < \lambda$, shrinks the bias $wak$. It has a price when there
are supply shocks. With $u = u^* - a(\pi - \pi^e) + e$, the discretionary solution becomes $\pi = \pi^{\mathrm T} + wak + \dfrac{wa}{1 + wa^2}e$ and $u = u^* + \dfrac{1}{1 + wa^2}e$: a bank with a small $w$ lets unemployment absorb almost
all of a shock. The best $w$ trades the bias against stabilisation, and lies strictly between zero and $\lambda$ (the weekend problem).

\begin{definition}[Inflation targeting]\label{def:b2:expectations-and-lucas-critique:inflation-targeting}
\emph{Inflation targeting}\index{inflation targeting} is a monetary-policy framework in which the central bank announces
a numerical target for inflation, is held accountable for meeting it, and explains its decisions in terms of the
forecast of inflation.
\end{definition}

The third way is a public target with accountability, which makes the cost of a surprise visible.
% ledger: b2:boe2012-count, b2:it-adopt-nzl
Twenty-seven countries ran a fully fledged inflation-targeting regime at the start of 2009, nine industrialised and eighteen
emerging and developing; New Zealand was first, in December 1989.

\begin{omchart}{Median CPI inflation in the 27 inflation-targeting countries, from five years before to ten years after the
  year of adoption (year 0), per cent.}{World Bank, World Development Indicators (FP.CPI.TOTL.ZG); adoption years from the
  Bank of England (2012), State of the art of \omterm{def:b2:expectations-and-lucas-critique:inflation-targeting}{inflation targeting}; our computation}
\begin{tikzpicture}
\begin{axis}[omrecord, width=0.6\linewidth, height=0.4\linewidth, xmin=-5.5, xmax=10.5, ymin=0, ymax=11,
    xlabel={Years from adoption}, ylabel={Median inflation, per cent}, xtick={-5,0,5,10}]
  \draw[black!40] (axis cs:0,0) -- (axis cs:0,11);
  \addplot[omDef, thick, mark=*, mark size=1.4pt] table[x=k, y=median]
    {figdata/bachelor-2/25-expectations-and-lucas-critique/out/it_event.csv};
\end{axis}
\end{tikzpicture}
\end{omchart}

% figdata: it_summary (k=-5 8.56, k=-1 6.06, k=0 6.37, k=+5 3.51, k=+10 3.03)
Median inflation in these countries was 8.56 per cent five years before adoption, 6.37 per cent in the year of adoption,
3.51 per cent five years after and 3.03 per cent ten years after. Most of the fall came after the target was set, but
inflation was already falling before, and the 1990s and 2000s were years of falling inflation in many countries that did
not adopt targets. The chart shows what happened in the targeters, not what would have happened without targets.

\begin{debate}[Does central-bank independence lower inflation?]
% ledger: b2:as1993-cbi-esp, b2:as1993-cbi-deu, b2:as1993-cbi-che, b2:as1993-abstract, b2:campillo1997
Alesina and Summers (1993), for 16 rich countries over 1955--1988, found that more independent central banks went with lower
and less variable inflation and with no measurable difference in real growth or unemployment: Germany and Switzerland,
at the top of their index, averaged 3.0 and 3.2 per cent inflation, Spain and New Zealand, at the bottom, 8.5 and 7.6.
Our averages in the opening chart, from annual consumer prices, are a little higher than the authors' (Spain 9.95 against
8.5, Germany 3.19 against 3.0), and New Zealand is not in our database, so the chart has 15 countries.
Campillo and Miron (1997), with a wider set of countries and controls, found that institutional arrangements such as
independence are relatively unimportant determinants of inflation, and economic fundamentals such as openness relatively
important. Two difficulties remain open. A legal index may differ from practice, and a country whose voters dislike
inflation may both choose an independent bank and keep inflation low for other reasons. How much of the correlation in the
opening chart is the effect of independence is not settled.
\end{debate}

\section{Survey measures of expectations}\label{sec:b2:expectations-and-lucas-critique:surveys}

\begin{measure}[How expectations are surveyed]
% ledger: b2:spf
Expectations cannot be observed in prices alone, so they are asked for. The oldest quarterly survey of macroeconomic
forecasts in the United States, the Survey of Professional Forecasters, began in 1968 and has been run by the Federal
Reserve Bank of Philadelphia since 1990; each quarter a panel of professional forecasters, identified only by number, give
their forecasts of output, unemployment and prices, and of inflation over the next year and the next ten years. Surveys of
households ask a sample of people how much they expect prices to rise. A survey measure records what respondents say;
whether expectations are rational is then a testable question: are the survey's forecast errors centred on zero, and
unrelated to what was known when the forecast was made?
\end{measure}

Surveys also measure credibility. When long-horizon expected inflation stays close to the target while current inflation
moves, people expect the bank to bring inflation back: the target is \omterm{def:b2:expectations-and-lucas-critique:credibility}{credible}. When long-horizon expectations follow
current inflation, it is not.

\section{Exercises}\label{sec:b2:expectations-and-lucas-critique:exercises}

\begin{exercise}[$\star$]\label{exo:b2:expectations-and-lucas-critique:1}
Read the opening chart. Which country had the highest average inflation, and which the lowest? State the slope of the
fitted line in words.
\end{exercise}

\begin{exercise}[$\star$]\label{exo:b2:expectations-and-lucas-critique:2}
Read the chart of American inflation and unemployment. Compare the 1960s, the 1970s and the 1980s.
\end{exercise}

\begin{exercise}[$\star$]\label{exo:b2:expectations-and-lucas-critique:3}
Read the inflation-targeting chart: median inflation five years before, at, and ten years after adoption. Why does the
chart not show the effect of targeting?
\end{exercise}

\begin{exercise}[$\star$]\label{exo:b2:expectations-and-lucas-critique:4}
Inflation has risen by one point a year for five years. Under adaptive expectations $\pi^e_t = \pi_{t-1}$, what is the forecast error each year? What
would \omterm{def:b2:expectations-and-lucas-critique:rational-expectations}{rational expectations} require of these errors?
\end{exercise}

\begin{exercise}[$\star\star$]\label{exo:b2:expectations-and-lucas-critique:5}
With $\pi^{\mathrm T} = 2$, $a = 0.5$, $k = 2$, $w = 1$ and no shocks, compute inflation, unemployment and the loss under discretion and under commitment. Find the bias if
$k$ falls to 1.
\end{exercise}

\begin{exercise}[$\star\star$]\label{exo:b2:expectations-and-lucas-critique:6}
Predict, then compute: in the listing's model, the target rises from 2 to 4 per cent. By how much does the intercept of the
re-estimated Phillips curve rise in theory? Compare with the listing's estimate.
\end{exercise}

\begin{exercise}[$\star\star$]\label{exo:b2:expectations-and-lucas-critique:7}
In the Lucas supply model, the bank uses the rule $\pi_t = 2 - 0.5(u_{t-1} - 5) + \varepsilon_t$ with $a = 2$. Last year unemployment was 7. What is expected inflation this year? If $\varepsilon_t = 1$, what are
inflation and unemployment?
\end{exercise}

\begin{exercise}[$\star\star$]\label{exo:b2:expectations-and-lucas-critique:8}
Modify the listing so that people expect last year's inflation instead of the target. What does the estimate now predict,
and is it right?
\end{exercise}

\begin{exercise}[$\star\star$]\label{exo:b2:expectations-and-lucas-critique:9}
\textbf{(D)}
With \texttt{econdata.jst}, compute average annual CPI inflation by decade, 1950--2019, for each of the eighteen countries. Report the
median across countries for each decade, the highest country average in the 1970s and 1980s, and which country it was.
\end{exercise}

\begin{exercise}[$\star\star\star$]\label{exo:b2:expectations-and-lucas-critique:10}
Derive \eqref{eq:b2:expectations-and-lucas-critique:reaction} and the discretionary equilibrium. Show that the bias is larger when surprises are more effective ($a$ larger).
\end{exercise}

\begin{exercise}[$\star\star\star$]\label{exo:b2:expectations-and-lucas-critique:11}
With the supply shock $e$, derive the discretionary inflation and unemployment given in
\cref{sec:b2:expectations-and-lucas-critique:credibility}.
\end{exercise}

\begin{exercise}[$\star\star\star$]\label{exo:b2:expectations-and-lucas-critique:12}
A government can commit to a state-contingent rule $\pi = \pi^{\mathrm T} + ce$. Find the $c$ that minimises society's expected loss, and show that this rule does
better than discretion with $w = \lambda$ by exactly the loss from the bias.
\end{exercise}

\section{Problem: Brenland's central-bank law}\label{sec:b2:expectations-and-lucas-critique:problem}

\begin{problem}[{Weekend problem --- Brenland's central-bank law: the loss function, discretion, a conservative central
banker, and the stabilisation trade-off}]
\label{pb:b2:expectations-and-lucas-critique:1}
Brenland's unemployment is $u = 5 - 0.5(\pi - \pi^e) + e$, with $e$ of mean zero and variance $\sigma^2 = 4$. Society's loss is $\tfrac12(\pi - 2)^2 + \tfrac12(u - 3)^2$. Parliament is drafting a law that will
give monetary policy to a bank with loss $\tfrac12(\pi - 2)^2 + \tfrac12w(u - 3)^2$.

\medskip
\noindent\textbf{Part I --- The loss function.}
\begin{enumerate}
  \item Identify $u^*$, $a$, $\pi^{\mathrm T}$, $k$ and $\lambda$.
  \item What does $k > 0$ say about Brenland's labour market?
  \item Why is the loss quadratic rather than linear?
  \item Without shocks, what is society's loss at $\pi = 2$ and $u = 5$?
  \item Can any policy do better than this on average, with \omterm{def:b2:expectations-and-lucas-critique:rational-expectations}{rational expectations} and no shocks? Why?
\end{enumerate}

\medskip
\noindent\textbf{Part II --- Discretion.} Set $w = 1$ and, for now, $e = 0$.
\begin{enumerate}[resume]
  \item Write the bank's choice of $\pi$ as a function of $\pi^e$.
  \item Find inflation in the equilibrium with \omterm{def:b2:expectations-and-lucas-critique:rational-expectations}{rational expectations}.
  \item Find unemployment.
  \item Compute society's loss.
  \item Compute the loss under commitment to 2 per cent.
  \item Why would an announcement of 2 per cent not be believed?
\end{enumerate}

\medskip
\noindent\textbf{Part III --- A conservative central banker.} Shocks now occur.
\begin{enumerate}[resume]
  \item Find the \omterm{def:b2:expectations-and-lucas-critique:inflation-bias}{inflation bias} as a function of $w$, and its value for $w = 0.5$.
  \item Find the responses of inflation and unemployment to a shock $e$ for $w = 1$ and for $w = 0$.
  \item Compute society's expected loss under discretion with $w = 1$.
  \item Compute it under the best commitment, with no bias and the responses of $w = 1$.
  \item Compute it with $w = 0$, a bank that cares only about inflation. Compare with $w = 1$.
  \item Compute it with $w = 0.5$.
\end{enumerate}

\medskip
\noindent\textbf{Part IV --- The stabilisation trade-off.}
\begin{enumerate}[resume]
  \item Write society's expected loss as a function of $w$.
  \item Find, numerically, the $w$ that minimises it.
  \item What is the expected loss at that $w$, and how far is it from the best commitment?
  \item What are the bias and the responses to a shock at that $w$?
  \item Why does the best $w$ lie strictly between 0 and 1?
  \item How does the best $w$ change if shocks are four times as variable? If $k$ doubles instead?
  \item State the result: the \omterm{def:b2:expectations-and-lucas-critique:inflation-bias}{inflation bias} under discretion and the weight that minimises expected social loss.
\end{enumerate}
\end{problem}
